\subsection{The complete cubic abelian calculation and the precise depth maps}

In this subsection $E=X_3^+$ is the cubic
$x^3+y^3+z^3=0$, $D=E\setminus U_3$, and
$G=J_{E,D}$ is the generalized Jacobian for the reduced divisor $D$.
We use the origin $O=[1:-1:0]$, and set
$P=[0:-1:1]$, $Q=[-1:0:1]$.
For an abelian group $A$ we write
$A^\wedge=\varprojlim_m A/mA$, with $m$ ordered by divisibility.
This is the completion for multiplication by $m$; the quotients
need not be finite. All inverse limits of cohomology groups below
carry their inverse-limit prodiscrete topology.
For a countable set $\mathcal I$ define
\[
 \mathcal R_{\mathcal I}(\widehat{\mathbb Z})
 =\left\{(a_j)\in\prod_{j\in\mathcal I}\widehat{\mathbb Z}:
       \text{for every }m,\ a_j\in m\widehat{\mathbb Z}
       \text{ for all but finitely many }j\right\}.
\]
Its topology is the one coming from
$\varprojlim_m\bigoplus_{\mathcal I}\mathbb Z/m$.
In particular this space is not asserted to be compact.

\begin{statement}{E.24}{Proved}{All cubic abelian sections and every fiber over the toric minimum}
Put $K=\mathbb Q(\zeta)$, $\zeta^2+\zeta+1=0$,
and $\lambda=1-\zeta$. There is a retraction of algebraic groups
over $\mathbb Q$
\[
 r:G\longrightarrow\mathbb G_m^2,\qquad
 G\simeq\mathbb G_m^2\times G',
 \qquad
 1\longrightarrow(\operatorname{Res}_{K/\mathbb Q}\mathbb G_m)^3
 \longrightarrow G'\longrightarrow E\longrightarrow0.
\]
The map $r$ is the map induced by the two units $x/z,y/z$ on $U_3$.
The group $E(\mathbb Q)$ is cyclic of order three, and
\[
 G'(\mathbb Q)\simeq\mu_6^3\oplus
                  \bigoplus_{j\in\mathcal I}\mathbb Z .
\]
Here one may take $\mathcal I$ to consist of three copies of the
nonzero prime ideals of $\mathbb Z[\zeta]$, with the first copy of
$(\lambda)$ replaced by one new generator. The displayed abstract
isomorphism includes explicit choices of prime generators.

Choose the unit rational tangent at $O$, with local parameter $z/x$,
as the origin of the set of sections of the geometrically abelian
arithmetic fundamental group. Up to geometric conjugation, that
entire set is then the group
\[
 H^1(G_{\mathbb Q},T_\bullet G)
 \simeq
 (\mathbb Q^{\times,\wedge})^2
 \times\mu_6^3
 \times\mathcal R_{\mathcal I}(\widehat{\mathbb Z})
 \times\widehat{\mathbb Z}.
 \tag{E24.1}
\]
The last factor and its splitting are noncanonical. The actual map
to the toric minimum of E.16 is the projection onto the first two
factors followed by translation by the Kummer class of $(1,-1)$.
Consequently it is surjective, and every fiber is a torsor under
\[
 \mu_6^3\times
 \mathcal R_{\mathcal I}(\widehat{\mathbb Z})\times\widehat{\mathbb Z}.
 \tag{E24.2}
\]
For the pro-$3$ geometric abelianization the same formulas hold
with $\mathbb Q^{\times,\wedge}$ replaced by
$\varprojlim_a\mathbb Q^\times/\mathbb Q^{\times3^a}$,
$\mu_6^3$ replaced by $\mu_3^3$, and each
$\widehat{\mathbb Z}$ replaced by $\mathbb Z_3$.
The translation by $-1$ is zero in this pro-$3$ version.
\end{statement}
\begin{proof}
\emph{The rational elliptic-curve calculation.}
The maps
\[
 X=-\frac{3z}{x+y},\qquad Y=\frac{4x-5y}{x+y},
 \qquad
 [x:y:z]=[Y+5:4-Y:-3X]
\]
identify the smooth cubic with $Y^2+Y=X^3-7$, with $O$ at infinity.
Indeed the exact identity is
\[
 Y^2+Y-X^3+7
 =\frac{27(x^3+y^3+z^3)}{(x+y)^3}.
\]
The inverse formula proves the birational assertion, which extends
to an isomorphism of the smooth projective curves.

For the rank calculation we use the explicit identification over
$\mathbb Q$ of the Fermat cubic with $X_0(27)$ in
Li--Long--Tu, \S2.4 and Remark 4.2
\cite{EC_LiLongTu}. The corresponding normalized differential has
newform
$f(\tau)=\eta(3\tau)^2\eta(9\tau)^2$.
The Mellin formula therefore gives
\[
 L(E,1)=2\pi\int_0^\infty
           \eta(3it)^2\eta(9it)^2\,dt>0 .
\]
The product formula for $\eta$ makes the integrand strictly
positive for $t>0$, and cuspidality gives convergence at both ends.
The automorphism induced by multiplication of one Fermat coordinate
by $\zeta$ gives complex multiplication by
$\mathbb Z[\zeta]$; this is the maximal order of the norm-Euclidean
field $K$. Coates--Wiles, Theorem 1 on p.~223
\cite{EC_CoatesWiles}, now gives rank zero.
The smooth cubic has exactly three points over $\mathbb F_2$,
so the prime-to-$2$ rational torsion injects into a group of order
three. An order-two point on the displayed Weierstrass model would
have $Y=-1/2$ and $X^3=27/4$, which has no rational solution.
Thus there is no $2$-primary torsion. The distinct points $O,P,Q$
already give three rational points, so $E(\mathbb Q)=\langle P\rangle
\simeq\mathbb Z/3$.

\emph{The retraction and the exact rational-point presentation.}
The cusp scheme is the disjoint union of three rational points and
three copies of $\operatorname{Spec}K$. Taking the frame at $O$ as
one identifies the torus of $G$ with
\[
 T_D=\mathbb G_m^2
       \times(\operatorname{Res}_{K/\mathbb Q}\mathbb G_m)^3.
\]
We write its coordinates as $(a,b,A,B,C)$, at $P,Q$ and the three
quadratic cusp orbits, respectively. The convention is that a frame
is a nonzero vector in the fiber of the line bundle.

Here is an explicit construction of the characters used in $r$.
For a unit $h$ on $U_3$, write
$\operatorname{div}(h)=\sum_{d\in D}m_d d$.
For a framed line bundle $(L,(e_d))$, choose a rational section
$s$ nonvanishing along $D$, and set
\[
 \chi_h(L,(e_d))
 =h(\operatorname{div}s)\,
        \prod_{d\in D}(s(d)/e_d)^{-m_d}.
 \tag{E24.3}
\]
Products over nonrational points mean the corresponding norms.
Replacing $s$ by $fs$ leaves this expression unchanged, since
Weil reciprocity gives
$h(\operatorname{div}f)=f(\operatorname{div}h)$.
Tensor products multiply (E24.3). The construction in families, or
the divisor presentation of the generalized Jacobian, consequently
gives a morphism of algebraic groups. On the torus it is
$\prod e_d^{m_d}$, and on a divisor $u-u_0$ in $U_3$ it is
$h(u)/h(u_0)$. Thus it is precisely the character induced by $h$.

For $h=x/z,y/z$ their divisors are the respective coordinate-cusp
orbits minus the orbit at infinity. Hence
\[
 r_T(a,b,A,B,C)
  =\bigl(aN(A)/N(C),\ bN(B)/N(C)\bigr).
 \tag{E24.4}
\]
The map $i(a,b)=(a,b,1,1,1)$ is an algebraic section of $r$.
The isomorphism with the direct product is explicitly
$g\mapsto(r(g),\,g\,i(r(g))^{-1})$.
The kernel torus is parametrized by
\[
 (A,B,C)\longmapsto
 \bigl(N(C)/N(A),\,N(C)/N(B),\,A,B,C\bigr).
 \tag{E24.5}
\]
Hilbert 90 gives $H^1(\mathbb Q,T_D)=0$ and
$H^1(\mathbb Q,(\operatorname{Res}_{K/\mathbb Q}\mathbb G_m)^3)=0$.

To calculate the extension on rational points, let $z_0\in G(\mathbb Q)$
be the framed line bundle $\mathcal O(P-O)$ with the local frames
specified next. Put
$f=(y+z)/(x+y)$, so
$\operatorname{div}(f)=3P-3O$.
At $P$ use $s=x/z$ and the frame $s^{-1}$ for
$\mathcal O(P-O)$; at $O$ use $s=z/x$ and the frame $s$.
At the other cusps use the canonical section $1$.
Multiplication by $f$ trivializes its third tensor power.
The two leading values of $f$ are
\[
 \lim_{P}f/s^3=1/3,\qquad
 \lim_{O}fs^3=3,\qquad f(Q)=-1 .
\]
For example $y/z=-(1+s^3)^{1/3}$ at $P$, which proves the
first equality; the same expansion at $O$ proves the second.
At the nonrational cusps use the representatives
$[0:-\zeta:1]$, $[-\zeta:0:1]$, $[1:-\zeta:0]$.
Dividing every frame value by the value $3$ at $O$ gives
\[
 3z_0=
 \left(3^{-2},-3^{-1},
             \lambda^{-1},\lambda^{-2},\zeta^2\lambda^{-3}\right)
       \in\mathbb Q^{\times2}\times K^{\times3}.
 \tag{E24.6}
\]
This calculation also follows directly from
$3=-\zeta^2\lambda^2$.
Equation (E24.4), with $N(\lambda)=3$, sends (E24.6) to $(1,-1)$.
The rational cube roots of $1$ and $-1$ are unique, so
$r(z_0)=(1,-1)$.
For $z'=z_0-i(1,-1)$ the presentation becomes
\[
 G'(\mathbb Q)=
 (K^{\times3}\oplus\mathbb Z z')/
 \left\langle 3z'-
   (\lambda^{-1},\lambda^{-2},\zeta^2\lambda^{-3})\right\rangle .
 \tag{E24.7}
\]
There are no additional relations: the quotient onto
$E(\mathbb Q)=\mathbb Z/3$ is surjective by Hilbert 90, the class of
$z'$ generates it, and (E24.7) specifies its third multiple in the
kernel.

The norm-Euclidean ring $\mathbb Z[\zeta]$ has units $\mu_6$,
so prime factorization gives
$K^\times=\mu_6\oplus\bigoplus_{\mathfrak p}\mathbb Z$ after choosing
one generator for each prime ideal. In (E24.7) the coefficient of
the first $\lambda$ generator is $-1$. Eliminate that generator
and retain $z'$ instead. This operation changes no torsion relation
and proves the asserted description of $G'(\mathbb Q)$.
As an independent finite-support check,
\texttt{e\_cubic\_abelian.py} obtains the Smith diagonal
$1,6,6,6$ for the relation involving the three $\mu_6$ generators,
the three $\lambda$ generators and $z'$.

\emph{The full cohomology calculation.}
Write $B=(\operatorname{Res}_{K/\mathbb Q}\mathbb G_m)^3$.
The extension of $E$ by $B$ in $G'$ is killed by three: the
extension of $E$ by $T_D$ is classified by the cusp-difference
map, all of whose values are killed by three by C.5, and (E24.5)
and the direct product exhibit the component defining $G'$.
It follows that the connecting map
\[
 \delta:H^1(\mathbb Q,E)\longrightarrow H^2(\mathbb Q,B)
\]
is killed by three. Since $H^1(\mathbb Q,B)=0$,
$H^1(\mathbb Q,G')$ injects into $H^1(\mathbb Q,E)$.
Every element of the maximal divisible subgroup of the latter
is three times an element of that subgroup; hence its image under
$\delta$ is zero. Consequently the injection identifies maximal
divisible subgroups:
\[
 \operatorname{Div}H^1(\mathbb Q,G')
   =\operatorname{Div}H^1(\mathbb Q,E)\simeq\mathbb Q/\mathbb Z .
 \tag{E24.8}
\]
The last equality is exactly
\c{C}iperiani--Stix, Theorem 18(2), printed p.~18
\cite{EC_CiperianiStix}, applied to the analytic-rank-zero curve
already established above. Their additional restrictions on primes
concern the separate assertion about the Tate--Shafarevich
intersection. They impose no restriction on (E24.8).

The Kummer sequences and inverse limit give
\[
 0\longrightarrow G'(\mathbb Q)^\wedge
 \longrightarrow H^1(G_{\mathbb Q},T_\bullet G')
 \longrightarrow
 \operatorname{Hom}(\mathbb Q/\mathbb Z,H^1(\mathbb Q,G'))
 \longrightarrow0 .
 \tag{E24.9}
\]
For precision, take the countable cofinal subsystem $m=a!$.
The transitions $G'(\mathbb Q)/m\to G'(\mathbb Q)/d$, $d\mid m$,
are surjective, so their first derived inverse limit vanishes.
The finite groups $H^0(\mathbb Q,G'[m])$ satisfy the
Mittag--Leffler condition, which identifies continuous $H^1$ of
$T_\bullet G'$ with the inverse limit of the finite-coefficient
$H^1$ groups. No finiteness of
$G'(\mathbb Q)/m$ or $H^1(\mathbb Q,G'[m])$ is used.
The last term in (E24.9) is the limit of
$H^1(\mathbb Q,G')[m]$ under multiplication transitions.
Every homomorphic image of $\mathbb Q/\mathbb Z$ is divisible,
so (E24.8) identifies it with
$\operatorname{End}(\mathbb Q/\mathbb Z)=\widehat{\mathbb Z}$.

Choose a lift of $1$ in (E24.9). Multiplication by
$\widehat{\mathbb Z}$ then gives a continuous splitting: modulo
$m$ its value depends only on the scalar modulo $m$.
The description (E24.7) gives
\[
 G'(\mathbb Q)^\wedge
  =\mu_6^3\times\mathcal R_{\mathcal I}(\widehat{\mathbb Z}).
 \tag{E24.10}
\]
Indeed the image at each $m$ has finite support in $\mathcal I$,
and compatibility is exactly the restricted-sequence condition
in the definition. This proves both the algebraic and
topological assertions in (E24.9).
The identical argument with the subsystem $m=3^a$ gives the
pro-$3$ version, whose divisible quotient is
$\operatorname{Hom}(\mathbb Q_3/\mathbb Z_3,
 \mathbb Q_3/\mathbb Z_3)=\mathbb Z_3$.

Finally, $T_\bullet G=H_1(U_{3,\overline{\mathbb Q}},
\widehat{\mathbb Z})$ by the generalized-Jacobian realization in
C.1. A rational tangential section splits the corresponding
arithmetic extension; its conjugacy classes of splittings are
exactly continuous one-cocycles modulo coboundaries.
Thus (E24.9), (E24.10), Hilbert 90 and the algebraic retraction
give (E24.1). At $O$, with $s=z/x$, the leading terms of
$(x/z,y/z)$ are $(s^{-1},-s^{-1})$. This makes the toric origin
the Kummer class of $(1,-1)$ and proves the claimed affine map.
The covering pushforward $q_{3*}$ is multiplication by three
after $r_*$ on geometric homology; the map used here is the
actual coordinate-unit map $r_*$, as in E.16. This fixes the
comparison and the translation, and proves the fiber assertion.
\end{proof}

\begin{statement}{E.25}{Proved}{Intrinsic depth, inherited depth, and the even-level obstruction}
Let $N_n=\pi_1(U_{n,\overline{\mathbb Q}})$, embedded in the
geometric profinite group $\widehat F$ of $U_1$ by the covering.
In this statement write $F=\widehat F$ in the profinite formulas.
Use $y$ for
inertia at the base puncture $1$.
For an intrinsic depth filtration choose a Galois-stable set of
cusps above $1$, and let $I_n$ be the closed normal subgroup
generated by their inertia. Put $D_{\mathrm{intr}}^r=\gamma_r I_n$.
One allowed choice for $n=3$ is the rational cusp $Q$ alone;
for $n=4$ take the entire four-point coordinate-cusp orbit.
There are actual quotient maps
\[
 N_n/D_{\mathrm{intr}}^3
   \longrightarrow
 \operatorname{im}(N_n\longrightarrow F/D_y^3),
 \qquad
 \ker=\frac{N_n\cap D_y^3}{D_{\mathrm{intr}}^3},
 \tag{E25.1}
\]
where $D_y^3=\gamma_3\langle\!\langle y\rangle\!\rangle$.
The same formulas apply to the corresponding Lie ideals and
their completed rational Lie algebras.

Over each field $\mathbb Q_\ell$ the inherited prounipotent
image is all of $F_{\mathbb Q_\ell}^{\mathrm{un}}/D_y^3$.
The corresponding inherited arithmetic extension is the base
arithmetic extension and has a rational tangential section,
for every $n$. On abelianized Lie algebras the map from the
intrinsic quotient has a kernel of dimension $n^2-1$.

For every even $n$, the inherited pro-$2$ depth quotient,
the inherited pro-$2$ class-two quotient, and every intrinsic
quotient retaining the curve's own pro-$2$ abelianization have
empty section sets over $\mathbb Q$.
The obstruction is the explicit central extension with class
$\kappa_2(-1)\cup\kappa_2(-1)$ described below.
\end{statement}
\begin{proof}
Inertia above $1$ maps to a conjugate of a power of $y$.
Consequently $I_n$ maps into $\langle\!\langle y\rangle\!\rangle$,
and lower central series give (E25.1), including its exact kernel.
The chosen sets of cusps are Galois stable, so their normal
closures and lower central terms are preserved by the full
arithmetic fundamental group. These are therefore arithmetic
quotients, not just quotients of unmarked geometric groups.

By E.9 the map
$N_{n,\mathbb Q_\ell}^{\mathrm{un}}\to
 F_{\mathbb Q_\ell}^{\mathrm{un}}$ is surjective: positive powers
of $x,y$ lie in $N_n$, and unique roots in characteristic-zero
unipotent groups recover $x,y$. Its inherited quotient is thus
the entire base quotient. The map on arithmetic extensions
induces this surjection on geometric kernels and the identity
on $G_{\mathbb Q}$. Quotienting by the geometric kernel of that
surjection identifies the inherited arithmetic extension with
the base one, which is split by the unit tangent at $0$ towards
$1$. Equivalently its fibers are the full fibers of
$(F_{\mathbb Q_\ell}^{\mathrm{un}}/D_y^3)\rtimes G_{\mathbb Q}$.
Since both depth ideals lie in the respective derived Lie
algebras, their abelianizations are the original ones.
The kernel dimension is therefore exactly the $n^2-1$ of B.5
and E.13.

The integral obstruction uses the quotient
\[
 F\longrightarrow
 D_8=\langle x,y,c\mid
   x^2=y^2=c^2=1,\ c\text{ central},\ [x,y]=c\rangle .
 \tag{E25.2}
\]
With the standard rational tangential section, the Galois action
on this quotient is trivial: in the normalized action of E.20
the cyclotomic exponent is odd, while the conjugating element
lies in $F'$ and hence has central image in $D_8$.
The arithmetic base quotient is $D_8\times G_{\mathbb Q}$.

In the plus model, with $t=-x^n/z^n$, the two covering roots are
\[
 (x/z)^n=-t,\qquad(y/z)^n=-(1-t).
 \tag{E25.3}
\]
At the unit tangent $t=0$, their constant root factors are
$n$th roots of $-1$. For even $n$, their reductions to quadratic
root factors both give
$\epsilon=\kappa_2(-1):G_{\mathbb Q}\to\mathbb Z/2$.
The sign in the stabilizer condition has no effect modulo two.
The geometric image of $N_n$ in $D_8$ is exactly
$\langle c\rangle$: its two exponent sums are even, and
$[x,y]\in N_n$ maps to $c$. The arithmetic image is consequently
the explicit pullback
\[
 \mathcal P=
 \{(d,g)\in D_8\times G_{\mathbb Q}:
           d^{\mathrm{ab}}=(\epsilon(g),\epsilon(g))\}.
 \tag{E25.4}
\]
For each $g$ there is a lift in $\pi_1(U_n)$, and multiplication
by the geometric element $c$ gives both lifts in (E25.4);
this proves equality of the image, not just its containment.
Using $(xy)^{\epsilon(g)}$ as a set-theoretic lift gives the
factor set $c^{\epsilon(g)\epsilon(h)}$.
Thus (E25.4) represents
$\kappa_2(-1)\cup\kappa_2(-1)$.

Complex conjugation has $\epsilon=1$. Both possible lifts,
$xy$ and $xyc$, square to $c\ne1$.
The extension cannot have a section even after restriction to
$G_{\mathbb R}$, hence cannot have one over $\mathbb Q$.
The group $D_8$ kills $\gamma_3F$ and hence $D_y^3$.
Its geometric image from $N_n$ is abelian, so it also kills
$N_n'$. These two explicit factorizations prove all three
nonexistence assertions in the statement.
They recover E.18 at level four and prove the assertion for
all even levels. Pushing to the rank-two toric lattice kills
the central element $c$, exactly as in E.18; extending scalars
to $\mathbb Q_\ell$ gives the inherited comparison already proved.
Thus neither operation is silently identified with the integral
intrinsic quotient.

For clarity the intrinsic marked-depth quotient also has a
precise comparison with the third derived quotient. Their common
quotient is
$N_n/(N_n'''D_{\mathrm{intr}}^3)$.
For the markings just specified neither kernel contains the other.
Indeed filling the marked cusps gives free groups of ranks
$9$ and $13$, respectively. On six distinct remaining generators
use the word
$w=[[[a,b],[a,c]],[[d,e],[d,f]]]$.
Its degree-eight Magnus term has the monomial $abacdedf$ with
coefficient one, and is nonzero. This third-derived word survives
the filling map,
so $N_n'''$ is not contained in $D_{\mathrm{intr}}^3$.
In the other direction choose a marked cusp generator $z_0$
and an unmarked free generator $a$. The element
$[z_0,[z_0,[z_0,a]]]$ lies in $\gamma_3 I_n$.
Its degree-four Magnus term is nonzero, whereas every element
of $N_n'''$ has zero Magnus terms below degree eight.
Thus it does not lie in $N_n'''$.
The corresponding maps from $N_n$, and the common quotient,
state the relation without imposing the base rank-two
containment from E.2 on a different free group.
\end{proof}
